\scanpage{138}
\mspage{9}
% archive sheet number: 137

\uline{Subdivisions}

Soient $A$ un \deleted{\unsure{organisme}} \inserted{atelier}, $\unsure{\mathcal{F}|\mathscr{F}}$ l'ensemble des $A$-figures, et
\[
G\ll F
\]
dans $\mathcal{F}$. On a donc une application \uline{croissante}
\diagram{\begin{tikzcd}[ampersand replacement=\&,row sep=small]
\unsure{\varphi|\phi}:\&\widetilde{G}\arrow[r]\arrow[d,equals]\&\widetilde{F}\arrow[d,equals]\\
\&G\&F
\end{tikzcd}}
telle que pour $Y\in\widetilde{G}$, $X=\varphi(Y)$, on ait
\[
Y\overset{\circ}{\ll}X
\]
Pour dire que $G$ est une \deleted{\unsure{raffinement}} \uline{subdivision} de $F$, nous exigerons \deleted{d'abord} que

(i) $\forall\,X\in\widetilde{F}$, \deleted{posant} considérant $\varphi^{-1}(X)\subset\widetilde{G}$,
$\varphi^{-1}(X)\overset{\circ}{\ll}X$, et $\operatorname{supp}^{\circ}\varphi^{-1}(X)\subset X^{\circ}$, on a
\deleted{on a}\inserted{donc}
\[
\deleted{\operatorname{supp}^{\circ}\varphi^{-1}(X)}
\]
\[
\operatorname{supp}^{\circ}\varphi^{-1}(X)=X^{\circ}
\]
% margin note written sideways
% slanted
\margin{Il valait mieux ici partir \unsure{d'une|d'un} \deleted{\unsure{subdivision|atelier}} \inserted{\unsure{figure}} \unsure{élém.} $F$ \unsure{donné|donnée} \unsure{et sa topologie}.}

En \unsure{admettant|admet.} \uline{\unsure{Atmag|Atomag} supp$^2$}, on voit que ceci équivaut à chacune des deux conditions

(ii)
\diagram{$\begin{array}{ccc}
\operatorname{supp}G&=&\operatorname{supp}F\\
\Vert&&\Vert\\
(\operatorname{supp}^{\unsure{\bullet|\circ}}G)&&(\operatorname{supp}^{\circ}F)
\end{array}$}
ou encore

(iii) Il n'existe pas $Z\in A$ \unsure{avec} $Z\ll F$, et $Z\mathrel{|\unsure{\circ|\bullet}|}\deleted{\unsure{F}}G$.
